\documentclass[12pt,a4paper]{article}
\usepackage[utf8]{inputenc}
\usepackage[T1]{fontenc}
\usepackage[polish]{babel}
\usepackage{geometry}
\usepackage{graphicx}
\usepackage{booktabs}
\usepackage{array}
\usepackage{longtable}
\usepackage{enumitem}
\usepackage{titlesec}
\usepackage{fancyhdr}
\usepackage{lastpage}
\usepackage{xcolor}
\usepackage{hyperref}

\geometry{margin=2.5cm}

% Kolory
\definecolor{pgblue}{RGB}{0,83,159}

% Nagłówek i stopka
\pagestyle{fancy}
\fancyhf{}
\fancyhead[L]{\small Dog FACS Dataset}
\fancyhead[R]{\small Raport semestralny I}
\fancyfoot[C]{\thepage\ / \pageref{LastPage}}
\renewcommand{\headrulewidth}{0.4pt}
\renewcommand{\footrulewidth}{0.4pt}

% Formatowanie sekcji
\titleformat{\section}{\large\bfseries\color{pgblue}}{\thesection.}{1em}{}
\titleformat{\subsection}{\normalsize\bfseries}{\thesubsection.}{1em}{}

\begin{document}

% Strona tytułowa
\begin{titlepage}
    \centering
    \vspace*{1cm}

    {\Large\textbf{Politechnika Gdańska}}\\[0.3cm]
    {\large Wydział Elektroniki, Telekomunikacji i Informatyki}\\[0.3cm]
    {\large Katedra Systemów Decyzyjnych i Robotyki}\\[2cm]

    {\Huge\textbf{Raport Semestralny}}\\[0.5cm]
    {\Large Semestr I (zimowy 2025/2026)}\\[2cm]

    {\LARGE\textbf{Dog FACS Dataset}}\\[0.3cm]
    {\large Utworzenie i Adnotacja Zbioru Danych\\do Analizy Emocji Psów}\\[3cm]

    \begin{tabular}{ll}
        \textbf{Zespół projektowy:} & Danylo Lohachov (koordynator)\\
        & Anton Shkrebela\\
        & Danylo Zherzdiev\\
        & Mariia Volkova\\[0.5cm]
        \textbf{Opiekun projektu:} & dr hab. inż. Michał Czubenko\\
    \end{tabular}

    \vfill

    {\large Gdańsk, styczeń 2026}
\end{titlepage}

\tableofcontents
\newpage

% ============================================================================
\section{Założenia projektu}
% ============================================================================

\subsection{Cel główny}

Celem projektu jest stworzenie wysokiej jakości, publicznie dostępnego zbioru danych do analizy emocji psów w formacie COCO, wspieranego przez dedykowany pipeline AI zdolny do automatycznej adnotacji.

\subsection{Założenia na semestr I}

Na semestr I zaplanowano następujące cele:

\begin{enumerate}[leftmargin=2cm]
    \item \textbf{Analiza i badania} -- analiza systemu DogFACS, formatu COCO, istniejących zbiorów danych i modeli SOTA
    \item \textbf{Model detekcji psów} -- fine-tuning YOLOv8 do wykrywania psów na obrazach (cel: mAP > 85\%)
    \item \textbf{Model klasyfikacji ras} -- fine-tuning EfficientNet do rozpoznawania ras psów (cel: Top-5 accuracy > 80\%)
    \item \textbf{Model punktów kluczowych} -- detekcja punktów kluczowych pyska psa (cel: PCK@0.1 > 75\%)
    \item \textbf{Klasyfikator emocji} -- mapowanie punktów kluczowych na emocje DogFACS (cel: accuracy > 70\%)
    \item \textbf{Pipeline inference} -- integracja wszystkich modeli w zunifikowany pipeline
    \item \textbf{Aplikacja demo} -- interfejs webowy (Streamlit) do wizualizacji wyników
    \item \textbf{Dokumentacja} -- DPP, DTP, dokumentacja techniczna
\end{enumerate}

\subsection{Metryki sukcesu}

\begin{table}[h]
\centering
\begin{tabular}{lcc}
\toprule
\textbf{Model} & \textbf{Metryka} & \textbf{Wymagane minimum} \\
\midrule
Detekcja BBox & mAP@0.5 & > 85\% \\
Klasyfikacja ras & Top-5 accuracy & > 80\% \\
Detekcja keypoints & PCK@0.1 & > 75\% \\
Klasyfikacja emocji & Accuracy & > 70\% \\
\bottomrule
\end{tabular}
\caption{Założone metryki jakości modeli AI}
\end{table}

% ============================================================================
\section{Zrealizowany zakres prac}
% ============================================================================

\subsection{Podsumowanie realizacji}

\textbf{Wszystkie główne założenia na semestr I zostały zrealizowane.} Poniżej przedstawiono szczegółowy status każdego zadania:

\begin{table}[h]
\centering
\begin{tabular}{p{6cm}ccc}
\toprule
\textbf{Zadanie} & \textbf{Status} & \textbf{Wynik} & \textbf{Cel} \\
\midrule
Analiza DogFACS i COCO & \textcolor{green!60!black}{Zrealizowane} & -- & -- \\
Model detekcji psów (YOLOv8m) & \textcolor{green!60!black}{Zrealizowane} & 87.3\% & > 85\% \\
Model klasyfikacji ras (EfficientNet-B4) & \textcolor{green!60!black}{Zrealizowane} & 82.1\% & > 80\% \\
Model keypoints (SimpleBaseline) & \textcolor{green!60!black}{Zrealizowane} & 76.8\% & > 75\% \\
Klasyfikator emocji (Rule-based) & \textcolor{green!60!black}{Zrealizowane} & 71.5\% & > 70\% \\
Pipeline inference & \textcolor{green!60!black}{Zrealizowane} & -- & -- \\
Aplikacja demo (Streamlit) & \textcolor{green!60!black}{Zrealizowane} & -- & -- \\
Dokumentacja DPP, DTP & \textcolor{green!60!black}{Zrealizowane} & -- & -- \\
\bottomrule
\end{tabular}
\caption{Status realizacji zadań semestru I}
\end{table}

\subsection{Produkty końcowe semestru I}

\begin{itemize}
    \item \textbf{Pipeline AI} -- działający system przetwarzający obrazy end-to-end:
    \begin{itemize}
        \item Detekcja psów → Klasyfikacja rasy → Detekcja keypoints → Klasyfikacja emocji
        \item Wyjście w formacie COCO JSON
    \end{itemize}
    \item \textbf{Aplikacja demo} -- interfejs Streamlit do wgrywania obrazów i wizualizacji wyników
    \item \textbf{Repozytorium GitHub} -- kod źródłowy z CI/CD
    \item \textbf{Dokumentacja} -- DPP, DTP, dokumentacja techniczna w LaTeX
\end{itemize}

% ============================================================================
\section{Zmiany względem pierwotnego planu}
% ============================================================================

W trakcie realizacji projektu wprowadzono dwie zmiany architektoniczne:

\subsection{Zmiana 1: Architektura modelu keypoints}

\begin{itemize}
    \item \textbf{Plan:} HRNet (High-Resolution Network)
    \item \textbf{Realizacja:} SimpleBaseline
    \item \textbf{Powód:} Lepsza dostępność pretrenowanych wag, prostsza implementacja przy porównywalnej jakości. HRNet wymagał znacznie więcej zasobów obliczeniowych i czasu na fine-tuning.
    \item \textbf{Wynik:} Model SimpleBaseline osiągnął PCK@0.1 = 76.8\%, przekraczając założony cel (> 75\%).
\end{itemize}

\subsection{Zmiana 2: Podejście do klasyfikacji emocji}

\begin{itemize}
    \item \textbf{Plan:} Trenowany model ML (sieć neuronowa)
    \item \textbf{Realizacja:} Rule-based classifier oparty na DogFACS
    \item \textbf{Powód:} Brak odpowiedniego zbioru treningowego z etykietami emocji psów. Istniejące datasety nie zawierają adnotacji emocji zgodnych z DogFACS.
    \item \textbf{Wynik:} Klasyfikator rule-based mapuje wykryte Action Units na kategorie emocji z accuracy = 71.5\%, przekraczając założony cel (> 70\%).
\end{itemize}

\textbf{Obie zmiany zostały zatwierdzone przez opiekuna projektu.}

% ============================================================================
\section{Wkład poszczególnych członków zespołu}
% ============================================================================

\subsection{Danylo Lohachov (U1) -- Koordynator projektu}

\begin{itemize}
    \item Koordynacja prac zespołu i harmonogramu
    \item Prowadzenie dokumentacji projektowej (DPP, DTP)
    \item Rozwój aplikacji demonstracyjnej (Streamlit)
    \item Zarządzanie repozytorium GitHub i Linear
    \item Kontrola jakości (code review)
\end{itemize}

\textbf{Główne osiągnięcia:}
\begin{itemize}
    \item Działająca aplikacja demo Streamlit
    \item Kompletna dokumentacja projektowa
    \item Skuteczna koordynacja 4-osobowego zespołu
\end{itemize}

\subsection{Anton Shkrebela (U2) -- Specjalista AI/ML}

\begin{itemize}
    \item Analiza i wybór architektury modelu keypoints
    \item Implementacja i trenowanie modelu detekcji punktów kluczowych
    \item Implementacja systemu DogFACS (11 Action Units)
    \item Opracowanie klasyfikatora emocji (rule-based)
    \item Ewaluacja i optymalizacja modeli
\end{itemize}

\textbf{Główne osiągnięcia:}
\begin{itemize}
    \item Model keypoints z PCK@0.1 = 76.8\%
    \item Działający klasyfikator emocji z accuracy = 71.5\%
    \item Kompletna implementacja systemu DogFACS
\end{itemize}

\subsection{Danylo Zherzdiev (U3) -- Backend Developer}

\begin{itemize}
    \item Fine-tuning modelu YOLOv8 do detekcji psów
    \item Fine-tuning modelu EfficientNet do klasyfikacji ras
    \item Integracja wszystkich modeli w zunifikowany pipeline
    \item Implementacja formatu COCO z rozszerzeniami DogFACS
    \item Optymalizacja wydajności pipeline'u
\end{itemize}

\textbf{Główne osiągnięcia:}
\begin{itemize}
    \item Model detekcji z mAP = 87.3\%
    \item Model klasyfikacji ras z Top-5 accuracy = 82.1\%
    \item Działający pipeline przetwarzający obrazy end-to-end
\end{itemize}

\subsection{Mariia Volkova (U4) -- Data Engineer}

\begin{itemize}
    \item Konfiguracja środowiska deweloperskiego
    \item Przygotowanie i przetwarzanie zbiorów danych treningowych
    \item Analiza istniejących datasetów (HuggingFace, Kaggle)
    \item Przygotowanie skryptów do pobierania i przetwarzania danych
    \item Wsparcie w testowaniu modeli
\end{itemize}

\textbf{Główne osiągnięcia:}
\begin{itemize}
    \item Skonfigurowane środowisko deweloperskie
    \item Przygotowane dane treningowe dla wszystkich modeli
    \item Skrypty do ekstrakcji klatek z filmów
\end{itemize}

\subsection{Podsumowanie wkładu}

\begin{table}[h]
\centering
\begin{tabular}{p{4cm}p{7cm}c}
\toprule
\textbf{Osoba} & \textbf{Główny obszar odpowiedzialności} & \textbf{Wkład} \\
\midrule
Danylo Lohachov & Koordynacja, dokumentacja, frontend & 30\% \\
Anton Shkrebela & Modele keypoints i emocji, DogFACS & 30\% \\
Danylo Zherzdiev & Modele BBox i ras, pipeline, COCO & 20\% \\
Mariia Volkova & Środowisko, dane, wsparcie & 20\% \\
\bottomrule
\end{tabular}
\caption{Podział wkładu w projekcie}
\end{table}

% ============================================================================
\section{Proponowane oceny}
% ============================================================================

Na podstawie zrealizowanego zakresu prac, osiągniętych wyników i zaangażowania poszczególnych członków zespołu, proponujemy następujące oceny:

\begin{table}[h]
\centering
\begin{tabular}{lcp{8cm}}
\toprule
\textbf{Osoba} & \textbf{Ocena} & \textbf{Uzasadnienie} \\
\midrule
Danylo Lohachov & 5.0 & Skuteczna koordynacja projektu, terminowa realizacja wszystkich zadań, kompletna dokumentacja, działająca aplikacja demo \\
\addlinespace
Anton Shkrebela & 5.0 & Samodzielna implementacja modeli keypoints i emocji, osiągnięcie założonych metryk, kreatywne rozwiązanie problemu braku danych treningowych \\
\addlinespace
Danylo Zherzdiev & 4.0 & Implementacja modeli detekcji i klasyfikacji, integracja pipeline'u \\
\addlinespace
Mariia Volkova & 4.0 & Konfiguracja środowiska, przygotowanie danych treningowych, wsparcie zespołu \\
\bottomrule
\end{tabular}
\caption{Proponowane oceny za semestr I}
\end{table}

\subsection{Uzasadnienie ocen}

\begin{enumerate}
    \item \textbf{Wszystkie założone cele zostały zrealizowane} -- każdy model AI osiągnął lub przekroczył wymagane metryki jakości
    \item \textbf{Terminowość} -- projekt realizowany zgodnie z harmonogramem, bez opóźnień
    \item \textbf{Jakość produktów} -- działający pipeline, aplikacja demo, kompletna dokumentacja
    \item \textbf{Współpraca zespołowa} -- równomierny podział pracy, skuteczna komunikacja, wspólne rozwiązywanie problemów
    \item \textbf{Adaptacja do zmian} -- elastyczne podejście do zmian architektonicznych przy zachowaniu celów projektu
\end{enumerate}

% ============================================================================
\section{Plany na semestr II}
% ============================================================================

W semestrze II planowane są następujące prace:

\begin{enumerate}
    \item \textbf{Zbieranie danych} -- pobranie około 2500 filmów psów z YouTube
    \item \textbf{Automatyczna adnotacja} -- przetworzenie minimum 25 000 klatek przez pipeline AI
    \item \textbf{Weryfikacja manualna} -- ręczna weryfikacja 25\% klatek (około 6250)
    \item \textbf{Klasyfikator emocji oparty na sieci neuronowej} -- zastąpienie rule-based classifiera modelem ML trenowanym na zebranych i zweryfikowanych danych z adnotacjami emocji
    \item \textbf{Rozwój aplikacji do ręcznej adnotacji} -- rozbudowa aplikacji Streamlit o pełny interfejs do manualnej adnotacji (bounding boxes, keypoints, emocje), umożliwiający weryfikację i korektę automatycznych adnotacji
    \item \textbf{Finalizacja datasetu} -- konwersja do formatu COCO, statystyki, dokumentacja
\end{enumerate}

\vspace{1cm}
\hrule
\vspace{0.5cm}

\textbf{Dokument przygotowany przez zespół projektowy Dog FACS Dataset}\\
Gdańsk, 29 stycznia 2026

\end{document}